% page 73 du fac-simile (image p-072 du scan)
% bloc 165.1 x 246.3 mm ; marge gauche 36.9 mm
\pgc{15.240mm}{0.000mm}{
\l{\cel{55}65}
\l{}
\l{}
\l{\sou{bilanco}. (\sou{komerco}).Equilibrigo di aktivo e di pasivo di konto}
\l{\cel{3}komercala. - DEFIRS.}
\l{}
\l{\sou{bilboketo}. Ludilo qua konsistas ye bastoneto qua finas, ca-}
\l{\cel{3}latere per pinto, talatere per pleto, mezi di qua esas}
\l{\cel{3}suspendita bulo quan onu devas lansar tale ke onu retro-}
\l{\cel{3}recevos lu sur la pinto o sur la pleto. - DFIR.}
\l{}
\l{\sou{bilgo}. Kolektuyo (en la fondo di la holdo di navo) di la}
\l{\cel{3}aqui, rezidui ed exkrementi, quan onu vakuigas per pumpili.-}
\l{\cel{3}DE.}
\l{}
\l{\sou{biljardo}. I. Pleo qua konsistas ye pulsar ivoro-buli per}
\l{\cel{3}longa bastono. - II. Tablo sur qua esas tensita tapiso verda,}
\l{\cel{3}garnisita ye rebordi, e sur qua onu pleas biliardo. - DEFIRS.}
\l{}
\l{\sou{bilieto}. Paper-folieto, karto, per qua konstatesas, favore di}
\l{\cel{3}ulu, yuro tempala, komprita, od obtenita per favoreso.-DEFIRS.}
\l{}
\l{\sou{biliono}. l,OOO,OOO2\cel{2}- DEFIRS.}
\l{}
\l{\sou{binara}. (\sou{pri sistemo di nombrifado)}. En qua onu expresas omna}
\l{\cel{3}nombri per du cifri (l e O). - DEFIS.}
\l{}
\l{\sou{bindar}. (\sou{trans}.)Sutar ad-ensemble la kayeri di libro, inkastrar}
\l{\cel{3}li en kovrilo kartona di qua la dorso e la platajo kovresas}
\l{\cel{3}da telo, felo, e c. - DE.}
\l{}
\l{\sou{bineto}. (\sou{agrokultivo}).Implemento kun mancho longa, di qua la}
\l{\cel{3}ferajo prizentas ye ca laItero di la mufo, lamo streta e}
\l{\cel{3}plata ; talatere, du (e mem tri) denti - e qua uzesas por}
\l{\cel{3}kultive modifikar duesmafoye sulo o planto. - FS.}
\l{}
\l{\sou{binoklo}. Du okulo-vitrodiski solidara. - DEFIRS.}
\l{}
\l{\sou{biodinamiko}. Cienco qua relatas la teorio (tre anciena) segun}
\l{\cel{3}qua existas vivo-forci qui diferas de la forci fiziko-kemiala,}
\l{\cel{3}e qui esus la fonto di la manifesti specifika : sentiveso,}
\l{\cel{3}moviveso, genetiveso, morto, e c., di la enti organoza.}
\l{\cel{3}- DEFIRS.}
\l{}
\l{\sou{binomio}.-(\sou{algebro}). Quanto qua konsistas ye du termini unionita}
\l{\cel{3}per la signo \textquotedbl{}plus\textquotedbl{} o la signo \textquotedbl{}minus\textquotedbl{}. - DEFIRS.}
\l{}
\l{\sou{biogenezo}. Teorio segun qua omna ento vivanta venis de ento}
\l{\cel{3}qua genitis lu. - DE.}
\l{}
\l{\sou{biografio}. Redakturo qua naracas la viv-eventaji di persono.}
\l{\cel{3}- DEFIRS.}
\l{}
\l{\sou{biokemio}. Parto di biologio, qua traktas la relati di la}
\l{\cel{3}strukturo kemiala di la elementi nemediata, a la mekanismo}
\l{\cel{3}funcionala di la celuli, tisui ed organi di la organismi}
\l{\cel{3}vivanta. - DEFIRS.}
}
